\appendix
\setcounter{equation}{0}
\setcounter{figure}{0}
\setcounter{table}{0}
\renewcommand{\theequation}{S\arabic{equation}}
\renewcommand{\thefigure}{S\arabic{figure}}
\renewcommand{\thetable}{S\arabic{table}}

\section*{Supplementary Information}

\subsection*{Supplementary Methods}

\subsubsection*{Parameter structure}

The punctuated Gamma model has four epidemiological parameters: the basic  reproduction number $R_0$, the early-epidemic growth rate $r$, and the two  parameters of the generation interval distribution, $\alpha$ (shape) and  $\beta$ (rate), which together determine the mean generation time  $\bar{g} = \alpha/\beta$. These four parameters are not all free: the  Euler--Lotka equation
%
\begin{equation}
	1  
	=R_0 \int_{\tau=0}^\infty e^{-r\tau } \frac{\beta^\alpha}{\Gamma(\alpha)} \tau^{\alpha-1} e^{-\beta \tau} d\tau
	= R_0 \left(\frac{\beta}{\beta + r}\right)^\alpha
	\label{eq:euler_lotka}
\end{equation}
%
provides one constraint among them, leaving three free parameters. We  choose to work with $(R_0, \bar{g}, \alpha)$ as the free parameters,  since each has a direct epidemiological interpretation. The rate parameter $\beta$ follows immediately from the  definition of the mean generation interval:
%
\begin{equation}
    \beta = \frac{\alpha}{\bar{g}},
\end{equation}
%
and the early-epidemic growth rate $r$ follows from the Euler--Lotka equation:
%
\begin{equation}
    r = \beta\left(R_0^{1/\alpha} - 1\right) 
      = \frac{\alpha}{\bar{g}}\left(R_0^{1/\alpha} - 1\right).
    \label{eq:r_from_params}
\end{equation}
%
The punctuation  parameter $\psi \in [0, 1]$ is the only remaining free parameter, and  it is the sole quantity that is not identifiable from population-level  epidemic data: since $A(\tau)$ is invariant across $\psi$ by  construction, the parameters $(R_0, \bar{g}, \alpha)$ can in principle  be estimated from the observed epidemic trajectory independently of  $\psi$. The consequences of varying $\psi$ while holding  $(R_0, \bar{g}, \alpha)$ fixed are therefore attributable purely to  changes in the individual-level temporal structure of infectiousness, with no confounding from changes in the population-level transmission  kernel.

\clearpage

\subsubsection*{Derivation of $E[\varsigma]$ and $\text{Var}[\varsigma]$}

{\bf Setup.} Following \citet{morris_computation_2024}, we describe the cumulative infections during the early-epidemic exponential growth phase as a branching process, $Z(t)$. Let $r$ be the epidemic's exponential growth rate; then, as $t \rightarrow \infty$, $e^{-rt} Z(t) \rightarrow W$ almost surely, where $W$ is a random variable that effectively scales the branching process' initial condition, capturing stochastic fluctuations in the early-epidemic process. We denote this convergence as $Z(t) \approx W \cdot e^{rt}$. 

% Here, $r$ is the early-epidemic exponential growth rate and $W$ is a random variable that effectively scales the branching process' initial condition, capturing stochastic fluctuations in the early-epidemic process. 

With a change of variables, we can translate $W$ into a random initial time shift: 
%
\begin{equation}
	Z(t) \approx e^{r (t + \varsigma)} \qquad \text{where} \qquad \varsigma = \frac{\log W - \log E[W]}{r} = \frac{\log W}{r}
\end{equation}
%
when $I_0 = E[W] = 1$ (\ie, when the epidemic is seeded with a single infectious person). Our goal is to determine how $E[\varsigma]$ and $\text{Var}[\varsigma]$ depend on the punctuation parameter $\psi$. To achieve this, we begin by examining the distribution of $W$. 

\noindent {\bf \boldmath Approximating the distribution of $W$.} First, we note that the distribution of $W$ consists of a point mass at 0 with mass $q$ (the probability of epidemic extinction), combined with a smooth density with positive support. The point mass translates into a time shift of $\varsigma = \log(0) / r = -\infty$, corresponding to an epidemic that never establishes. Thus, we restrict our attention to $W | \text{surv}$ and $\varsigma | \text{surv}$, that is, $W$ and $\varsigma$ conditioned on epidemic survival. 

To approximate the distribution of $W|\text{surv}$, we use a $\text{Gamma}(k, \lambda)$ distribution matched to the first two moments. This is a version of the five-moment-matching approach described by \cite{morris_computation_2024}, simplified for analytic tractability. Specifically, we let  
\begin{equation}
	k = \frac{E[W|\text{surv}]^2}{\text{Var}[W|\text{surv}]} \qquad \text{and} \qquad \lambda = \frac{E[W|\text{surv}]}{\text{Var}[W|\text{surv}]}
	\label{eq:k_and_lambda}
\end{equation}
Our task is now to compute $E[W|\text{surv}]$ and $\text{Var}[W|\text{surv}]$. 

\noindent {\bf \boldmath Deriving $E[W|\text{surv}]$.} We note that 
\begin{equation}
E[W] = E[W|\text{surv}] P(\text{surv}) + E[W|\text{extinct}] P(\text{extinct}) = E[W | \text{surv}] \cdot p
\label{eq:p_explanation}
\end{equation}
since $E[W|\text{extinct}] = 0$, and where $p = 1-q$ is the probability of epidemic survival. Thus,
\begin{equation}
\boxed{E[W|\text{surv}] = \frac{1}{p}}
\end{equation}
\begin{center}
{\it [for plugging into expressions for $k$ and $\lambda$, Eq.~\ref{eq:k_and_lambda}]} \\ 
\end{center} 
under our assumption that $E[W] = 1$. 

\noindent {\bf \boldmath Deriving $\text{Var}[W|\text{surv}]$.} The derivation for $\text{Var}[W]$ is somewhat more involved. Recall the Gamma time-shift model, where an infectious person produces $\chi_i \sim \text{Poisson}(R_0)$ secondary infections distributed at times $\tau_j = l_i+ \epsilon_j,\, j\in \{1, \dots, \chi_i\}$ and 
\begin{equation}
l_i \sim \text{Gamma}((1-\psi) \alpha, \beta), \qquad \epsilon_j \sim \text{Gamma}(\psi \alpha, \beta)
\end{equation}
This ensures that a random secondary infection time $\tau$ is distributed according to the generation interval distribution $g(\tau)$, \ie,
\begin{equation}
	\tau \sim \text{Gamma}(\alpha, \beta)
\end{equation}
For notational convenience, we introduce 
\begin{equation}
	E[e^{-c r \tau}] = \int_0^\infty e^{-cr\tau} \frac{\beta^\alpha}{\Gamma(\alpha)} \tau^{\alpha-1} e^{-\beta \tau}\, d\tau
	 = \Bigl(\frac{1}{1 + cz}\Bigr)^\alpha := \rho_c^\alpha
\end{equation}
that is, 
\begin{equation}
\boxed{\rho_c = \frac{1}{1+cz}}
\end{equation}
\begin{center}
{\it [for notational convenience]} \\ 
\end{center}
where $z = r/\beta = r \bar{g} / \alpha = R_0^{1/\alpha} - 1$, a dimensionless quantity that simplifies expressions and allows one to easily plug in different combinations of epidemiological parameters. Here, $\bar{g}$ is the mean of the generation interval distribution $g(\tau)$, and $c$ is a positive integer. The final expression, $z = R_0^{1/\alpha} - 1$, comes from the Euler-Lotka equation \citep{wallinga_how_2007}: 
\begin{equation}
1 = R_0 E[e^{-r \tau}] = R_0 \rho_1^\alpha 
= R_0 \Bigl(\frac{1}{1+z}\Bigr)^\alpha
\label{eq:euler_lotka}
\end{equation}
% \begin{equation}
% 1 = R_0 E[e^{-r \tau}] = R_0 \int_0^\infty e^{-r \tau} \frac{\beta^\alpha}{\Gamma(\alpha)} \tau^{-\alpha-1} e^{-\beta \tau}\, d\tau
% = R_0 \Bigl(\frac{1}{1+z}\Bigr)^\alpha
% \end{equation}
and solving for $z$. 


Next, we introduce the distributional fixed point equation, which is a standard starting point for computing moments of $W$. This equation exploits the self-similarity of the branching process, in which an epidemic initiated by a single individual is, in distribution, a scaled and time-shifted copy of the whole epidemic. At large $t$, the epidemic seeded by the index case $i$ consists of the downstream epidemics caused by their offspring $j$, which have grown approximately to size $Z_j(t) = W_j e^{r(t - \tau_j)}$. That is, 
%
\[
	Z(t) \approx \sum_{j=1}^{\chi_i} W_j e^{rt} e^{-r \tau_j}
\]
Multiplying through by $e^{-rt}$: 
%
\[  W \overset{d}{=} \sum_{j=1}^{\chi_i} e^{-r \tau_j} W_j \]

We now use this expression to calculate $\text{Var}[W]$. Since randomness enters through both $W$ and $\tau$, we first condition on the number of offspring and their infection times, $\tau_\bullet = \{ \chi_i, \tau_1, \tau_2, \dots, \tau_{\chi_i} \}$, using the law of total variance: 
%
\[ \text{Var}[W] = E[\text{Var}\Bigl[\sum_j e^{-r \tau_j} W_j | \tau_\bullet\Bigr]] + \text{Var}[E\Bigl[\sum_j e^{-r \tau_j} W_j | \tau_\bullet\Bigr]] \]
\[
=E\Bigl[ \sum_j e^{-2r\tau_j}  \cdot \text{Var}[W] \Bigr] + 
\text{Var} \Bigl[\sum_j e^{-r \tau_j} \cdot E[W] \Bigr]
\]
\[
= E\Bigl[\sum_j e^{-2 r \tau_j}\Bigr] \cdot \text{Var}[W] + 
(E[W])^2 \cdot \text{Var}\Bigl[\sum_j e^{-r \tau_j}\Bigr]
\]
%
Since $E[W] = I_0 = 1$, we have 
\[
	\text{Var}[W] = \text{Var}[W] E\Bigl[\sum_j e^{-2 r \tau_j}\Bigr] + \text{Var}\Bigl[\sum_j e^{-r \tau_j}\Bigr]
\]
Solving for $\text{Var}[W]$, 
\[
	\text{Var}[W] = \frac{\text{Var}[\sum_j e^{-r \tau_j}]}{1 - E[\sum_j e^{-2 r \tau_j}]}
\]
To simplify the denominator, note that 
\[ 
E\Bigl[\sum_{j=1}^{\chi_i} e^{-2 r \tau_j}\Bigr] = E[\chi_i] E[e^{-2 r \tau_j}]
 = R_0 E[e^{-2 r \tau_j}]
 % = R_0 \Bigl( \frac{\beta}{\beta + 2r}\Bigr )^\alpha 
 = R_0 \rho_2^\alpha 
 \]
Thus, the variance expression simplifies to 
\[
	\text{Var}[W] = \frac{\text{Var}[\sum_j e^{-r \tau_j}]}{1 - R_0 \rho_2^\alpha}
\]
Next, we examine the numerator: 
\[
	\text{Var}[\sum_j e^{-r \tau_j}] = E[(\sum_j e^{-r \tau_j})^2] - E[\sum_j e^{-r \tau_j}]^2 = E[(\sum_j e^{-r \tau_j})^2] - 1
\]
where the last step follows from the Euler-Lotka equation (Eq.~\ref{eq:euler_lotka}). 
% The second term is (R_0 \frac{1}{1+z})
% \[
% 	\text{Var}[\sum_j e^{-r \tau_j}] = E[(\sum_j e^{-r \tau_j})^2] - E[\sum_j e^{-r \tau_j}]^2 = E[(\sum_j e^{- r \tau_j})^2] - 1
% \]
Then, 
\[
	E[(\sum_j e^{- r \tau_j})^2] = E[\sum_j e^{-2 r \tau_j}] + E[\sum_{j \neq k} e^{-r (\tau_j + \tau_k)}]
\]
\[= R_0 \rho_2^\alpha + E[\chi_i (\chi_i - 1)] E[e^{-r \tau_j} e^{-r \tau_k}]\]
\[= R_0 \rho_2^\alpha + R_0^2 E[e^{-r \tau_j} e^{-r \tau_k}]\]
Here, the $\chi_i (\chi_i - 1)$ term enters because there are this many terms in the sum over $j \neq k$. Also, 
\[
E[\chi_i (\chi_i - 1)] = E[\chi_i^2] - E[\chi_i] = \text{Var}[\chi_i] + E[\chi_i]^2 - E[\chi_i] = R_0 + R_0^2 - R_0 = R_0^2
\]
under the assumption that $\chi_i \sim \text{Poisson}(R_0)$. 

Since $\tau_j = l_i + \epsilon_j$, and since $l_i$, $\epsilon_j$, and $\epsilon_k$ are independent, 
\[ E[e^{-r \tau_j} e^{-r \tau_k}]  = E[e^{-2r l_i} e^{-r \epsilon_j} e^{-r \epsilon_k}] = E[e^{-2r l_i}] (E[e^{-r \epsilon}])^2 = \rho_2^{(1-\psi)\alpha} \cdot \rho _1^{2 \psi \alpha}\]
Finally, we can assemble $\text{Var}[W]$: 
\[
	\text{Var}[\sum_j e^{-r \tau_j}] = R_0 \rho_2^\alpha + R_0^2 \rho_2^{(1-\psi)\alpha} \rho_1^{2 \psi \alpha} - 1
\]
and thus, substituting in the $z$ expressions for $\rho_1$ and $\rho_2$ and noting $R_0 = (1 + z)^\alpha$, we obtain
\begin{equation}
	\text{Var}[W] = \frac{
		\Bigl(\frac{1 + z}{1 + 2z}\Bigr)^\alpha + \Bigl(\frac{(1 + z)^2}{1 + 2z}\Bigr)^{(1-\psi) \alpha} - 1
	}{
		1 - \Bigl(\frac{1 + z}{1 + 2z}\Bigr)^\alpha
	} := 
	\frac{\sigma_1^\alpha + \sigma_2^{(1-\psi)\alpha} - 1}{1 - \sigma_1^\alpha}
\end{equation}
where $\sigma_c = \frac{(1 + z)^c}{1 + 2z}$ for positive integer $c$. 

  Last, we note that $E[W^n|\text{surv}] = E[W^n]/p$
  (by the same logic as in Eq.~\ref{eq:p_explanation}), so
  \begin{align}
      \text{Var}[W | \text{surv}]
      &= E[W^2|\text{surv}] - E[W|\text{surv}]^2 = \frac{1}{p} E[W^2] - \frac{1}{p^2} E[W]^2 \\ 
      &= \boxed{\frac{1}{p} (1 + \text{Var}[W]) - \frac{1}{p^2}}
  \end{align}  
% Last, we note that, by the same logic as in Eq.~\ref{eq:p_explanation}, 
% \begin{equation}
% 	\text{Var}[W | \text{surv}] = \frac{1}{p} \cdot \frac{
% 		\Bigl(\frac{1 + z}{1 + 2z}\Bigr)^\alpha + \Bigl(\frac{(1 + z)^2}{1 + 2z}\Bigr)^{(1-\psi) \alpha} - 1
% 	}{
% 		1 - \Bigl(\frac{1 + z}{1 + 2z}\Bigr)^\alpha
% 	}
% \end{equation}
\begin{center}
{\it [for plugging into expressions for $k$ and $\lambda$, Eq.~\ref{eq:k_and_lambda}]}
\end{center}

\noindent {\bf \boldmath Expressions for $k$ and $\lambda$.} Combining the previous derivations and plugging into Eq.~\ref{eq:k_and_lambda} gives 
\begin{equation}
\boxed{k = \frac{1 - \sigma_1^\alpha}{p \sigma_2^{(1 - \psi)\alpha} - 1 + \sigma_1^\alpha},
\qquad
\lambda = p k }
\end{equation}
\begin{center}
	{\it [Parameters of the moment-matched Gamma distribution for W]}
\end{center}

\noindent {\bf \boldmath Deriving the approximate density of $\varsigma$.} Let $f_W(x)$ be the Gamma-distributed approximate density of $W$: 
\begin{equation}
	f_W(x) = \frac{\lambda^k}{\Gamma(k)} x^{k-1} e^{-\lambda x}
\end{equation}
Since $\varsigma = \frac{1}{r} \log W$, we have $W = e^{r \varsigma}$, or, in terms of the arguments of the density functions, $x = e^{r s}$. By the change of variables formula, the density of $\varsigma$ is 
\begin{equation}
f_\varsigma(s) = f_W(e^{rs}) |\frac{dw}{ds}| = 
\frac{\lambda^k}{\Gamma(k)} e^{rs(k-1)} e^{-\lambda e^{rs}} r e^{rs} = \frac{r \lambda^k}{\Gamma(k)} e^{rsk} e^{-\lambda e^{rs}}
\label{eq:varsigma_dist}
\end{equation}
By taking the logarithm, differentiating with respect to $s$, and setting equal to 0, we can determine the mode of the distribution: 
\begin{equation}
	\boxed{\text{Mode}[\varsigma] = \frac{1}{r} \log \Bigl(\frac{k}{\lambda}\Bigr) = \frac{1}{r} \log \frac{1}{p}}
	\label{eq:varsigma_mode}
\end{equation}
\begin{center}
	{\it [Mode of the time shift distribution]}
\end{center}
Thus, the modal time shift depends on the epidemic growth rate $r$ and the establishment probability $p$ (and thus the reproduction number $R_0$), but not on the punctuation parameter $\psi$. 


\noindent {\bf \boldmath Deriving $E[\varsigma]$.} 
It would be possible to derive $E[\varsigma]$ directly from Eq.\ref{eq:varsigma_dist}, but there is a more direct route: since $W \sim \text{Gamma}(k,\lambda)$ and $\varsigma = \frac{1}{r} \log(W)$, it follows that 
\begin{equation}
	E[\varsigma] = \frac{1}{r}(\Psi^{(0)}(k) - \log \lambda)
\end{equation}
where $\Psi^{(0)}(k)$ is the digamma function evaluated at $k$. Plugging in $\lambda = pk$: 
\begin{equation}
\boxed{E[\varsigma] = \frac{\log(1/p)}{r} + \frac{\Psi^{(0)}(k) - \log(k)}{r}}
\end{equation}
\begin{center}
	{\it [Expected value of the time shift distribution]}
\end{center}
The first term is $\text{Mode}[\varsigma]$ (Eq.~\ref{eq:varsigma_mode}), and is independent of the punctuation parameter $\psi$. The second term is a $\psi$-dependent adjustment that is: 
\begin{enumerate}
	\item {\bf  Strictly negative.} {\it Proof:} Let $G \sim \text{Gamma}(k,1)$. Then, $E[G] = k$ and $E[\log G] = \Psi^{(0)}(k)$. By Jensen's inequality, since $\log$ is strictly concave, 
	\[\Psi^{(0)}(k) = E[\log G] < \log(E[G]) = \log k \]
	Thus, $\bigl(\Psi^{(0)}(k) - \log(k)\bigr)/r < 0$. 
	\item {\bf \boldmath Strictly increasing (less negative) in $\psi$.} {\it Proof:} Consider 
	\[\frac{d}{dk} \bigl[ \Psi^{(0)}(k) - \log k \bigr] = 
	\frac{d}{dk} \Biggl[ \int_0^\infty \Bigl(\frac{e^{-t}}{t} - \frac{e^{-kt}}{1 - e^{-t}}\Bigr) dt \Biggr] - \frac{1}{k} 
	= \int_0^\infty \frac{t e^{-k t }}{1 - e^{-t}}\, dt  - \frac{1}{k}
	\]
	Note that $1 - e^{-t} < t$ for $t > 0$, and thus $\frac{t}{1 - e^{-t}} > 1$. So, 
	\[  
	\int_0^\infty \frac{t e^{-kt}}{1 - e^{-t}} > \int_0^\infty e^{-kt} = \frac{1}{k}
	 \]
	 So, 
	 \[\frac{d}{dk} \bigl[ \Psi^{(0)}(k) - \log k \bigr] > 0 \]
	 and thus $(\Psi^{(0)}(k) - \log k)/r$ is increasing in $k$. \\ 
	 Next, we note that $k$ is increasing in $\psi$. The only $\psi$-dependent quantity in $k$ (Eq.\ref{eq:k_and_lambda}) is $\sigma_2^{(1-\psi)\alpha}$, which shows up in the denominator. Since
	 \[\sigma_2 = \frac{(1 + z)^2}{(1 + 2z)} = 1 + \frac{z^2}{1 + 2z} > 1 \quad \text{ for  } z > 0\]
	 the term is maximized when $\psi = 0$ and minimized when $\psi = 1$. Thus, $k$ is minimized when $\psi = 0$ and maximized when $\psi = 1$ --- \ie, $k$ is increasing in $\psi$. 
\end{enumerate}

\noindent Together, these observations demonstrate that the expected time shift sits later than the modal time shift (\ie, the time shift is left-skewed; Observation 1), and the expected time shift moves increasingly later as $a_i(\tau)$ becomes more punctuated (Observation 2). 

\noindent {\bf \boldmath Deriving $\text{Var}[\varsigma]$.} Following again from the properties of $\log W$ where $W$ is $\text{Gamma}(k, \lambda)$-distributed, 
\begin{equation}
	\boxed{\text{Var}[\varsigma] = \frac{1}{r^2} \Psi^{(1)}(k)}
\end{equation}
\begin{center}
	{\it [Variance of the time shift distribution]}
\end{center}
where $\Psi^{(1)}(k)$ is the trigamma function. The trigamma function is strictly decreasing in $k$, since 
\begin{equation}
\frac{d}{dk} \Psi^{(1)}(k) =  \frac{d}{dk }\sum_{n=0}^\infty \frac{1}{(k+n)^2} = 
-2 \sum_{n=0}^\infty \frac{1}{(k + n)^3}
\end{equation}
which is strictly negative when $k > 0$. Thus, $\psi \rightarrow 0$ (punctuated profiles) $\Rightarrow$ $k$ decreasing $\Rightarrow$ $\text{Var}[\varsigma]$ increasing, \ie, more punctuated profiles increase the variance of the time shift. 

\noindent {\bf Summary.} In this section, we derived properties of the random time shift $\varsigma$ by relying on a branching process approximation to the early-epidemic exponential growth phase. When conditioning on epidemic survival, we demonstrated that the time shift becomes increasingly later in expectation and more variable as the individual infectiousness profile $a_i(\tau)$ becomes more punctuated ($\psi \rightarrow 0$).

\clearpage

\subsubsection*{Generality beyond the time-shifted Gamma model}

The preceding derivation used the time-shifted Gamma model and a Gamma-distributed generation interval. Here, we show that the core result --- punctuated profiles increase the variance of early-epidemic timing --- holds for \emph{any} non-degenerate generation interval distribution (\ie, any distribution with positive variance, as opposed to a degenerate distribution concentrated on a single point). The argument rests only on the additive splitting $\tau_j = l_i + \varepsilon_j$ (shared shift plus independent jitter) and the Cauchy-Schwarz inequality.

\noindent {\bf Setup.} Let $g$ be an arbitrary generation interval distribution with Laplace transform $M_g(s) = E[e^{-s\tau}]$, so the Euler-Lotka equation gives $R_0 M_g(r) = 1$. Each index case's offspring share a common latent shift $l_i$ and receive independent jitter $\varepsilon_j$, so $\tau_j = l_i + \varepsilon_j$ with the marginal distribution of $\tau_j$ equal to $g$. The distributions of $l$ and $\varepsilon$ are arbitrary, subject only to $M_l \cdot M_\varepsilon = M_g$.

\noindent {\bf \boldmath Model-free $\text{Var}[W]$ formula.} The law-of-total-variance derivation proceeds identically to the Gamma case. The only splitting-dependent step is the cross-term for $j \neq k$:
\[
E[e^{-r \tau_j} e^{-r \tau_k}] = E[e^{-2rl}] (E[e^{-r\varepsilon}])^2 = M_l(2r) \cdot M_\varepsilon(r)^2
\]
Using $M_l(2r) = M_g(2r) / M_\varepsilon(2r)$, this becomes $M_g(2r) \cdot Q$, where
\begin{equation}
	Q := \frac{M_\varepsilon(r)^2}{M_\varepsilon(2r)}
	\label{eq:Q_def}
\end{equation}
Assembling the full variance formula (with $A := R_0 M_g(2r)$):
\begin{equation}
	\text{Var}[W] = \frac{A(1 + R_0 Q) - 1}{1 - A}
	\label{eq:var_W_general}
\end{equation}
Here, $A$ depends only on $g$ and not on the splitting, and $Q$ is the \emph{only} place the splitting enters. Since Eq.~\ref{eq:var_W_general} is affine in $Q$ with positive slope $R_0 A / (1-A)$, $\text{Var}[W]$ is strictly increasing in $Q$.

\noindent {\bf Three properties of $Q$.}

\begin{enumerate}
	\item $Q \leq 1$, with equality if and only if $\varepsilon$ is degenerate (maximum punctuation).

	{\it Proof:} By the Cauchy-Schwarz inequality applied to $X = e^{-r\varepsilon}$:
	\[M_\varepsilon(2r) = E[X^2] \geq (E[X])^2 = M_\varepsilon(r)^2\]
	so $Q = M_\varepsilon(r)^2 / M_\varepsilon(2r) \leq 1$. Equality holds if and only if $X$ is almost surely constant, \ie, $\varepsilon$ is degenerate. $\blacksquare$

	\item $Q$ is multiplicative under convolution: if $\varepsilon_2 \overset{d}{=} \varepsilon_1 + \delta$ with $\delta$ independent, then $Q_2 = Q_1 \cdot Q_\delta \leq Q_1$, with strict inequality if $\delta$ is non-degenerate.

	{\it Proof:} $M_{\varepsilon_2}(s) = M_{\varepsilon_1}(s) \cdot M_\delta(s)$, so $Q_2 = Q_1 \cdot M_\delta(r)^2 / M_\delta(2r) = Q_1 \cdot Q_\delta$, and $Q_\delta \leq 1$ by Property 1. $\blacksquare$

	\item The spike-vs-smooth gap in $\text{Var}[W]$ is strictly positive for any non-degenerate $g$.

	{\it Proof:} For the spike ($\varepsilon = 0$), $Q_\text{spike} = 1$. For the smooth ($\varepsilon = \tau$, $l = 0$), $Q_\text{smooth} = M_g(r)^2 / M_g(2r) = 1 / (R_0^2 M_g(2r)) = 1/(R_0 A)$. The gap is
	\[
	\text{Var}[W]_\text{spike} - \text{Var}[W]_\text{smooth} = \frac{R_0 A (1 - Q_\text{smooth})}{1 - A} = \frac{R_0 A - 1}{1 - A}
	\]
	Since $R_0 A = R_0^2 M_g(2r) \geq R_0^2 M_g(r)^2 = 1$ by the Cauchy-Schwarz inequality applied to $X = e^{-r\tau}$, with strict inequality for non-degenerate $g$, the gap is strictly positive. $\blacksquare$
\end{enumerate}

\noindent {\bf Monotonicity for parameterized families.} For a one-parameter family of splittings $\varepsilon_\psi$ with $\varepsilon_0 = 0$ (spike) and $\varepsilon_1 = \tau$ (smooth), if the family is ordered by convolution --- meaning $\varepsilon_{\psi_2} \overset{d}{=} \varepsilon_{\psi_1} + \delta$ for some independent $\delta \geq 0$ whenever $\psi_2 > \psi_1$ --- then $Q(\psi)$ is strictly decreasing and $\text{Var}[W](\psi)$ is strictly decreasing in $\psi$. The time-shifted Gamma model satisfies this ordering, since $\text{Gamma}(\psi_2 \alpha, \beta) \overset{d}{=} \text{Gamma}(\psi_1 \alpha, \beta) + \text{Gamma}((\psi_2 - \psi_1) \alpha, \beta)$, but any infinitely divisible generation interval distribution admits such a family.

\noindent {\bf Implications for the full chain.} Combined with the model-free conditioning step ($\text{Var}[W|\text{surv}] = (1 + \text{Var}[W])/p - 1/p^2$, which increases in $\text{Var}[W]$) and the Gamma-approximation results above, the full chain holds for general $g$:
\begin{center}
\begin{tabular}{ccccccc}
	More punctuation & $\Rightarrow$ & higher $Q$ & $\Rightarrow$ & higher $\text{Var}[W]$ & $\Rightarrow$ & lower $k$ \\
	& & {\scriptsize (Cauchy-Schwarz)} & & {\scriptsize (Eq.~\ref{eq:var_W_general})} & & {\scriptsize (Eq.~\ref{eq:k_and_lambda})} \\[6pt]
	& $\Rightarrow$ & more negative $E[\varsigma]$ & and & higher $\text{Var}[\varsigma]$ \\
	& & {\scriptsize (digamma monotonicity)} & & {\scriptsize (trigamma monotonicity)} \\
\end{tabular}
\end{center}
The first two links are model-free: the only assumption is that $g$ is non-degenerate. The final link --- from $\text{Var}[W]$ to properties of $\varsigma = (\log W)/r$ --- requires a distributional assumption, because moments of $W$ do not uniquely determine moments of $\log W$. The Gamma moment-matching approximation serves this role: given $E[W|\text{surv}]$ and $\text{Var}[W|\text{surv}]$, it provides a complete distribution from which $E[\varsigma]$ and $\text{Var}[\varsigma]$ can be computed in closed form via the digamma and trigamma functions. Without this (or an analogous) approximation, the qualitative direction is still supported by Jensen's inequality --- which guarantees $E[\varsigma|\text{surv}] < \log(1/p)/r$ for any distribution of $W|\text{surv}$, with the gap increasing when the distribution becomes more dispersed in the sense of convex order --- but the specific quantitative expressions would not apply.











% \noindent To complete the argument, we show that $k$ is increasing in $\psi$. Recall from Eq.~\ref{eq:k_and_lambda} that
% \[
% k = \frac{1 - \sigma_1^\alpha}{p\, \sigma_2^{(1-\psi)\alpha} + \sigma_1^\alpha - 1}
% \]
% The numerator does not depend on $\psi$. In the denominator, the only $\psi$-dependent term is $\sigma_2^{(1-\psi)\alpha}$, whose derivative with respect to $\psi$ is
% \[
% \frac{d}{d\psi} \sigma_2^{(1-\psi)\alpha} = -\alpha \log(\sigma_2) \cdot \sigma_2^{(1-\psi)\alpha} < 0
% \]
% since $\sigma_2 = (1+z)^2/(1+2z) = 1 + z^2/(1+2z) > 1$ for $z > 0$. Thus, the denominator is strictly decreasing in $\psi$, and $k$ (a positive constant divided by a decreasing positive function) is strictly increasing in $\psi$.

% \noindent {\bf Summary.} Combining the three results above: as $\psi$ decreases (more punctuated infectiousness), $k$ decreases, making the Jensen gap $(\Psi^{(0)}(k) - \log k)/r$ more negative, and thus $E[\varsigma]$ more negative --- meaning epidemics reach the establishment threshold earlier, on average.

% secondary infection time are distributed according to the generation interval distribution $g(\tau)$, \ie, $\tau_j \sim \text{Gamma}(\alpha, \beta)$


% For notational convenience, we introduce 
% \begin{equation}
% 	\rho_c = E[e^{-c r \tau}] = 
% \end{equation}





% Our first step is to approximate the distribution of $W$, conditional on epidemic survival, with a $\text{Gamma}(k, \lambda)$ distribution by matching the first two moments. Specifically, we set
% % \begin{equation}
% % 	k = \frac{}
% % \end{equation}


% Our first step is to approximate the distribution of $W$, which we denote $f_W(x)$, with an approximate distribution $\tilde{f}_W(x)$ that follows a $\text{Gamma}(k,\lambda)$ distribution. 





% \clearpage


% \subsubsection*{Derivation of $\text{Var}[W]$}

% For a branching process $Z(t)$ conditional on non-extinction, $e^{-rt} Z(t) \rightarrow W$ almost surely as $t \rightarrow \infty$. Here, $r$ is the early-epidemic exponential growth rate and $W$ is a random variable that effectively scales the branching process' initial condition, capturing stochastic fluctuations in the early-epidemic process. 

% With a change of variables, we can translate $W$ into a random initial time shift: 
% %
% \begin{equation}
% 	\varsigma = \frac{\log W - \log E[W]}{r} = \frac{\log W}{r}
% \end{equation}
% %
% when $I_0 = E[W] = 1$. Our goal is to understand how $\text{Var}[\varsigma]$ depends on the punctuation parameter $\psi$. Since $\varsigma$ and $W$ are monotonically related, we can derive an expression for $\text{Var}[W]$ and will interpret its implications for $\text{Var}[\varsigma]$ at the end. 

% To start, we consider an index case $i$ infected at time $t = 0$ who produces $\chi_i \sim \text{Poisson}(R_0)$ offspring. Under the Gamma time-shift model, the offspring $j$ from this index case are infected at times $\tau_j = l_i + \epsilon_j$, where $l_i$ is person $i$'s latent period and $\epsilon_j$ is the additional lag between $l_i$ and person $j$'s infection time, with 
% %
% \[l_i \sim \text{Gamma}((1-\psi)\alpha, \beta), \qquad \epsilon_j \sim \text{Gamma}(\psi \alpha, \beta) \]
% %
% This ensures that the population-level generation interval distribution is
% %
% \[ g(\tau) \sim \text{Gamma}(\alpha, \beta) \]
% %
% We can define some useful quantities that will arise later in the derivation. First, from the Euler-Lotka equation under $g(\tau)$ \citep{wallinga_how_2007}:, 
% %
% \[1 = R_0 \int_{\tau=0}^\infty e^{-r\tau } \frac{\beta^\alpha}{\Gamma(\alpha)} \tau^{\alpha-1} e^{-\beta \tau} d\tau\]
% %
% which simplifies to 
% %
% \[ 1 = R_0 \Bigl( \frac{\beta}{\beta+r }\Bigr )^\alpha = R_0 \Bigl(\frac{1}{1+z}\Bigr)^\alpha\]
% %
% where  $z = r/\beta = r \bar{g} / \alpha = R_0^{1/\alpha} - 1$, a dimensionless quantity that simplifies expressions and allows one to easily plug in different combinations of epidemiological parameters. Here, $\bar{g}$ is the mean of the generation interval distribution $g(\tau)$. 

% Next, we define 
% %
% \[ \rho_k =  (E[e^{-k r \tau}])^{1/\alpha} = \frac{\beta}{\beta + k r}  = \frac{1 }{1 + kz}\]
% This allows us to re-write the Euler-Lotka equation: 
% %
% \[ 1 = R_0 \rho_1^\alpha \]

% Next, we introduce the distributional fixed point equation, which is a standard starting point for computing moments of $W$. This equation exploits the self-similarity of the branching process, in which an epidemic initiated by a single individual is, in distribution, a scaled and time-shifted copy of the whole epidemic. At large $t$, the epidemic seeded by the index case $i$ consists of the downstream epidemics caused by their offspring $j$, which have grown approximately to size $Z_j(t) = W_j e^{r(t - \tau_j)}$. That is, 
% %
% \[
% 	Z(t) \approx \sum_{j=1}^{\chi_i} W_j e^{rt} e^{-r \tau_j}
% \]
% Multiplying through by $e^{-rt}$: 
% %
% \[  W \overset{d}{=} \sum_{j=1}^{\chi_i} e^{-r \tau_j} W_j \]

% We now use this expression to calculate $\text{Var}[W]$. Since randomness enters through both $W$ and $\tau$, we first condition on $\tau_j$ using the law of total variance: 
% %
% \[ \text{Var}[W] = E[\text{Var}\Bigl[\sum_j e^{-r \tau_j} W_j | \tau_j\Bigr]] + \text{Var}[E\Bigl[\sum_j e^{-r \tau_j} W_j | \tau_j\Bigr]] \]
% \[
% =E\Bigl[ \sum_j e^{-2r\tau_j}  \cdot \text{Var}[W] \Bigr] + 
% \text{Var} \Bigl[\sum_j e^{-r \tau_j} \cdot E[W] \Bigr]
% \]
% \[
% = E\Bigl[\sum_j e^{-2 r \tau_j}\Bigr] \cdot \text{Var}[W] + 
% (E[W])^2 \cdot \text{Var}\Bigl[\sum_j e^{-r \tau_j}\Bigr]
% \]
% %
% Since $E[W] = I_0 = 1$, we have 
% \[
% 	\text{Var}[W] = \text{Var}[W] E\Bigl[\sum_j e^{-2 r \tau_j}\Bigr] + \text{Var}\Bigl[\sum_j e^{-r \tau_j}\Bigr]
% \]
% Solving for $\text{Var}[W]$, 
% \[
% 	\text{Var}[W] = \frac{\text{Var}[\sum_j e^{-r \tau_j}]}{1 - E[\sum_j e^{-2 r \tau_j}]}
% \]
% To simplify the denominator, note that 
% \[ 
% E\Bigl[\sum_{j=1}^{\chi_i} e^{-2 r \tau_j}\Bigr] = E[\chi_i] E[e^{-2 r \tau_j}]
%  = R_0 E[e^{-2 r \tau_j}]
%  % = R_0 \Bigl( \frac{\beta}{\beta + 2r}\Bigr )^\alpha 
%  = R_0 \rho_2^\alpha 
%  \]
% Thus, the variance expression simplifies to 
% \[
% 	\text{Var}[W] = \frac{\text{Var}[\sum_j e^{-r \tau_j}]}{1 - R_0 \rho_2^\alpha}
% \]
% The denominator does not depend on $\psi$, so we proceed by examining the numerator: 
% \[
% 	\text{Var}[\sum_j e^{-r \tau_j}] = E[(\sum_j e^{-r \tau_j})^2] - E[\sum_j e^{-r \tau_j}]^2 = E[(\sum_j e^{- r \tau_j})^2] - 1
% \]
% Then, 
% \[
% 	E[(\sum_j e^{- r \tau_j})^2] = E[\sum_j e^{-2 r \tau_j}] + E[\sum_{j \neq k} e^{-r (\tau_j + \tau_k)}]
% \]
% \[= R_0 \rho_2^\alpha + E[\chi_i (\chi_i - 1)] E[e^{-r \tau_j} e^{-r \tau_k}]\]
% \[= R_0 \rho_2^\alpha + R_0^2 E[e^{-r \tau_j} e^{-r \tau_k}]\]
% Here, the $\chi_i (\chi_i - 1)$ term enters because there are this many terms in the sum over $j \neq k$. Also, 
% \[
% E[\chi_i (\chi_i - 1)] = E[\chi_i^2] - E[\chi_i] = Var[\chi_i] + E[\chi_i]^2 - E[\chi_i] = R_0 + R_0^2 - R_0 = R_0^2
% \]
% under the assumption that $\chi_i \sim \text{Poisson}(R_0)$. 

% Since $\tau_j = l_i + \epsilon_j$, and since $l_i$, $\epsilon_j$, and $\epsilon_k$ are independent, 
% \[ E[e^{-r \tau_j} e^{-r \tau_k}]  = E[e^{-2r l_i} e^{-r \epsilon_j} e^{-r \epsilon_k}] = E[e^{-2r l_i}] (E[e^{-r \epsilon}])^2 = \rho_2^{(1-\psi)\alpha} \cdot \rho _1^{2 \psi \alpha}\]
% Finally, we can assemble $\text{Var}[W]$: 
% \[
% 	\text{Var}[\sum_j e^{-r \tau_j}] = R_0 \rho_2^\alpha + R_0^2 \rho_2^{(1-\psi)\alpha} \rho_1^{2 \psi \alpha} - 1
% \]
% and thus, substituting in the $z$ expressions for $\rho_1$ and $\rho_2$, we obtain
% \begin{equation}
% 	\text{Var}[W] = \frac{
% 		\Bigl(\frac{1 + z}{1 + 2z}\Bigr)^\alpha + \Bigl(\frac{(1 + z)^2}{1 + 2z}\Bigr)^{(1-\psi) \alpha} - 1
% 	}{
% 		1 - \Bigl(\frac{1 + z}{1 + 2z}\Bigr)^\alpha
% 	}
% \end{equation}
% which is Eq.~\ref{eq:var_W}. 

% \clearpage 

% \subsubsection*{Mode of the two-moment Gamma approximation for the time-shift distribution}

% Following a simplified version of the moment-matching approach of  \citet{morris_computation_2024}, we approximate the distribution of $W$ conditional on survival by a Gamma distribution matched to its first two moments. The unconditional moments of $W$ are:
% %
% \begin{align}
%     E[W] &= 1 \\[6pt]
%     E[W^2] &= 
%     \text{Var}[W] + E[W]^2 = 
%     % \frac{\sigma^{(1-\psi)\alpha}}{1 - c^\alpha} = 
%     \frac{
% 		\Bigl(\frac{(1 + z)^2}{1 + 2z}\Bigr)^{(1-\psi) \alpha}
% 	}{
% 		1 - \Bigl(\frac{1 + z}{1 + 2z}\Bigr)^\alpha
% 	}
% \end{align}
% %
% Conditioning on survival with probability $p$: 
% %
% \begin{align}
%     E[W \mid \text{surv}] &= \frac{E[W]}{p} = \frac{1}{p} 
%     \label{eq:W_cond_mean} \\[6pt]
%     E[W^2 \mid \text{surv}] &= \frac{E[W^2]}{p} 
%     % = \frac{\sigma^{(1-\psi)\alpha}}{p(1 - c^\alpha)}
%     = \frac{
% 		\Bigl(\frac{(1 + z)^2}{1 + 2z}\Bigr)^{(1-\psi) \alpha}
% 	}{
% 		p \Bigl(1 - \Bigl(\frac{1 + z}{1 + 2z}\Bigr)^\alpha\Bigr)
% 	}
%     \label{eq:W_cond_second_moment}
% \end{align}
% %
% where the division by $p$ arises because 
% \begin{equation}
% E[W^n] = E[W^n | \text{surv}] P(\text{surv}) + E[W^n | \text{extinct}] P(\text{extinct}) = E[W^n | \text{surv}] \cdot p 
% \end{equation}
% since $W = 0$ when the epidemic fails to establish. Thus, $E[W^n | \text{surv}] = E[W^n] / p$. 

% Our goal is to approximate the distribution of $W$ with a $\text{Gamma}(k, \lambda)$ distribution, and specifically to estimate the mode of the distribution of $W | \text{surv}$ with the mode of the moment-matched $\text{Gamma}(k, \lambda)$ distribution. Note that the mode of a generic $\text{Gamma}(k,\lambda)$ distribution is $(k-1)/\lambda$, and that the distribution has a mean of $k/\lambda$ and a squared coefficient of variation of ($\text{CV}^2$) $1/k$. Thus, the parameters of the moment-matched Gamma distribution are 
% \begin{align}
% 	k &= \frac{1}{\text{CV}^2[W|\text{surv}]}  \\ 
% 	\lambda &= \frac{k}{E[W|\text{surv}]} = k p 
% \end{align}
% and the mode of the moment-matched distribution is 
% \begin{equation}
% 	\text{mode}[W|\text{surv}] \approx \frac{k - 1}{\lambda} = \frac{k-1}{k p } = \frac{1 - 1/k}{p} = \frac{1 - \text{CV}^2[W|\text{surv}]}{p}
% \end{equation}
% The expression for $\text{CV}^2$ is  
% \begin{align}
% 	\text{CV}^2(W|\text{surv})  &= \frac{\text{Var}(W|\text{surv})}{E[W|\text{surv}]^2} 
% 	= \frac{E[W^2 | \text{surv}] - E[W| \text{surv}]^2}{E[W|\text{surv}]^2} \\ 
% 	&= \frac{
% 		p \Bigl(\frac{(1 + z)^2}{1 + 2z}\Bigr)^{(1-\psi) \alpha}
% 	}{
% 		1 - \Bigl(\frac{1 + z}{1 + 2z}\Bigr)^\alpha
% 	} - 1
% 	\label{eq:cv2}
% \end{align}
% This moment-matched approximation is valid when $k \geq 1$ (\ie,\ $\text{CV}^2[W|\text{surv}] \leq 1$). When $k < 1$, the  approximating density is monotonically decreasing on $(0, \infty)$ with  no interior mode. 

% The time-shift $\varsigma = \log(W)/r$ is a monotone transformation of $W$, so its modal value is obtained by transforming  the mode of the Gamma approximation. Therefore:
% %
% \begin{equation}
%     \text{Mode}[\varsigma] 
%     = \frac{1}{r}\log\!\left(\frac{1-\text{CV}^2[W|\text{surv}]}{p}\right)
%     \label{eq:mode_varsigma_split}
% \end{equation}
% %
% Substituting $\text{CV}^2[W|\text{surv}]$ from Eq.~\eqref{eq:cv2} and expressing in terms of  $z = R_0^{1/\alpha} - 1$, this yields Eq.~\eqref{eq:mode_timeshift}.  

% \clearpage 

% \subsubsection*{Monotonicity in $R_0$}

% Differentiating Eq.~\eqref{eq:alpha_log_sigma} with respect to $R_0$,  holding $\alpha$ fixed, and using  $\frac{d}{dR_0}(2R_0^{1/\alpha}) = (2/\alpha)R_0^{1/\alpha - 1}$:
% %
% \begin{equation}
%     \frac{\partial(\alpha\log\sigma)}{\partial R_0} 
%     = \frac{2}{R_0} - \frac{2R_0^{1/\alpha - 1}}{2R_0^{1/\alpha} - 1}
%     = \frac{2}{R_0} \cdot \frac{R_0^{1/\alpha} - 1}{2R_0^{1/\alpha} - 1}
%     = \frac{2(\mu - 1)}{R_0(2\mu - 1)}
% \end{equation}
% %
% This is strictly positive for all $R_0 > 1$, since $\mu = R_0^{1/\alpha} > 1$  makes both the numerator $\mu - 1$ and denominator $2\mu - 1$ positive. The  derivative vanishes at $R_0 = 1$ ($\mu = 1$), confirming that at the epidemic  threshold there is no $\psi$-dependence in either $\text{Var}(W)$ or  $\text{Mode}[\varsigma]$.

% \clearpage 

% \subsubsection*{Monotonicity in $\alpha$}

% Differentiating Eq.~\eqref{eq:alpha_log_sigma} with respect to $\alpha$,  holding $R_0$ fixed. Write $g(\alpha) = \alpha\log(2R_0^{1/\alpha} - 1)$,  so that $\alpha\log\sigma = 2\log R_0 - g(\alpha)$ and it suffices to show  $\partial g/\partial\alpha > 0$. By the product rule, using  $\frac{\partial}{\partial\alpha}R_0^{1/\alpha}  = -R_0^{1/\alpha}\log R_0/\alpha^2$:
% %
% \begin{equation}
%     \frac{\partial g}{\partial\alpha} 
%     = \log(2\mu - 1) - \frac{2\mu\log\mu}{2\mu - 1}
% \end{equation}
% %
% where we have substituted $\mu = R_0^{1/\alpha}$ and $\log R_0 = \alpha\log\mu$.  To establish the sign, define  $h(\mu) = (2\mu-1)\log(2\mu-1) - 2\mu\log\mu$, so that the condition  $\partial g/\partial\alpha > 0$ is equivalent to $h(\mu) > 0$.  Differentiating:
% %
% \begin{equation}
%     h'(\mu) = 2\log(2\mu - 1) + 2 - 2\log\mu - 2 
%             = 2\log\!\left(2 - \frac{1}{\mu}\right)
% \end{equation}
% %
% For $\mu > 1$, we have $2 - 1/\mu > 1$, so $h'(\mu) > 0$. Since $h(1) = 0$  and $h'(\mu) > 0$ for all $\mu > 1$, it follows that $h(\mu) > 0$ for all  $\mu > 1$, giving $\partial g/\partial\alpha > 0$ and hence:
% %
% \begin{equation}
%     \frac{\partial(\alpha\log\sigma)}{\partial\alpha} < 0
% \end{equation}
% %
% So $\alpha\log\sigma$ is strictly decreasing in $\alpha$: more concentrated  generation intervals (larger $\alpha$, holding $R_0$ fixed) reduce the  impact of punctuation on both $\text{Var}(W)$ and $\text{Mode}[\varsigma]$.

% \clearpage 

% \subsection*{The modulating role of $\alpha \log \sigma$}

% The $\psi$-dependent term in both $\text{Var}(W)$ and $\text{Mode}[\varsigma]$  is $\sigma^{(1-\psi)\alpha}$, which can be written as  $e^{(1-\psi)\alpha\log\sigma}$. The quantity $\alpha\log\sigma$ therefore  acts as the fundamental modulator of $\psi$-dependence: larger  $\alpha\log\sigma$ means the base $\sigma > 1$ is raised to a larger  effective power, amplifying the spread in $\text{Var}(W)$ and  $\text{Mode}[\varsigma]$ between the spike ($\psi \to 0$) and smooth  ($\psi \to 1$) limits. We can express this quantity exactly in terms of  $R_0$ and $\alpha$ alone. Starting from $\sigma = \mu^2/(2\mu-1)$ with  $\mu = R_0^{1/\alpha}$:
% %
% \begin{equation}
%     \log\sigma = 2\log\mu - \log(2\mu - 1) 
%                = \frac{2\log R_0}{\alpha} - \log(2R_0^{1/\alpha} - 1)
% \end{equation}
% %
% and multiplying through by $\alpha$:
% %
% \begin{equation}
%     \alpha\log\sigma = 2\log R_0 - \alpha\log(2R_0^{1/\alpha} - 1).
%     \label{eq:alpha_log_sigma}
% \end{equation}
% %
% We now show that this expression is strictly increasing in $R_0$ and  strictly decreasing in $\alpha$, confirming the qualitative picture that  punctuation ($\psi < 1$) matters more when $R_0$ is large and the  generation interval is dispersed (small $\alpha$).

% \clearpage 

% \subsubsection*{Summary}

% The two partial derivatives establish that the modulating quantity  $\alpha\log\sigma$ --- and hence the sensitivity of $\text{Var}(W)$ and  $\text{Mode}[\varsigma]$ to $\psi$ --- is:
% %
% \begin{itemize}
%     \item \textbf{Increasing in $R_0$}: higher transmissibility amplifies  the effect of punctuation, because more offspring more densely sample  the individual infectiousness profile, inheriting more of its temporal  structure.
%     \item \textbf{Decreasing in $\alpha$}: more concentrated generation  intervals dampen the effect of punctuation, because there is less  temporal variation across the generation interval for the shared  latent period $l_i$ to exploit.
% \end{itemize}
% %
% Together, these results confirm that punctuated infectiousness  ($\psi < 1$) has the greatest impact on epidemic timing variability  when $R_0$ is large and the generation interval is dispersed ---  precisely the regime in which the individual infectiousness profile  has the most temporal structure for offspring timing to inherit.

% \clearpage

% \subsubsection*{Overdispersion arising from periodic contact patterns}

% Here, we derive the Negative Binomial overdispersion parameter $k$ \citep{lloyd-smith_superspreading_2005} of the secondary infection distribution under the periodic contact model (Eq.~\ref{eq:periodic_c}), which we restate here in a stationary form with a common random phase $\Phi \sim \text{Uniform}[0, 2\pi)$:
% %
% \begin{equation}
% 	c_i(t) = 1 - A \cos(\omega t + \Phi), \qquad A \in [0,1].
% 	\label{eq:periodic_c_supp}
% \end{equation}
% %
% Randomizing $\Phi$ makes $c_i(\cdot)$ stationary and is equivalent to assuming that the infection-time phase $t_i \bmod 2\pi/\omega$ is uniform across individuals, which holds in good approximation during sustained exponential growth.

% \paragraph{Lloyd--Smith bridge.} The number of secondary infections from person $i$ is $Z_i \sim \text{Poisson}(M_i)$ with $M_i = \int_0^\infty a_i(t_i, \tau)\,d\tau = \nu_i\,\tilde{c}_i$, where
% %
% \begin{equation}
% 	\tilde{c}_i := \int_0^\infty b_i(\tau)\,c_i(t_i + \tau)\,d\tau.
% 	\label{eq:tilde_c_def}
% \end{equation}
% %
% Matching the mean and variance of the mixed-Poisson offspring law to those of $\text{NegBin}(R_0, k)$ in the Lloyd-Smith parameterization gives $k = R_0^2/\text{Var}[M_i]$. With $\nu_i \equiv R_0$ held fixed (isolating the contribution of punctuation and contact variability), this reduces to
% %
% \begin{equation}
% 	k = \frac{1}{\text{Var}[\tilde{c}_i]}.
% 	\label{eq:k_bridge}
% \end{equation}
% %
% The rest of the derivation computes $\text{Var}[\tilde{c}_i]$.

% \paragraph{Mean.} Conditional on $b_i$, stationarity gives $E[c_i(t)] = 1$, hence
% %
% \begin{equation}
% 	E[\tilde{c}_i \mid b_i] = \int_0^\infty b_i(\tau)\,d\tau = 1,
% \end{equation}
% %
% using that $b_i$ is a probability density. Therefore $E[\tilde{c}_i] = 1$ (and $E[M_i] = R_0$), and by the law of total variance $\text{Var}[\tilde{c}_i] = E\bigl[\text{Var}[\tilde{c}_i \mid b_i]\bigr]$.

% \paragraph{Autocovariance.} Averaging over the uniform random phase $\Phi$,
% %
% \begin{equation}
% 	\text{Cov}[c_i(t), c_i(t+s)] = A^2\,E_\Phi[\cos(\omega t + \Phi)\cos(\omega(t+s) + \Phi)] = \frac{A^2}{2}\cos(\omega s),
% 	\label{eq:periodic_cov}
% \end{equation}
% %
% so $\text{Var}[c_i(t)] = A^2/2$ and the spectrum of the contact process is concentrated at the single frequency $\omega$.

% \paragraph{Conditional variance of $\tilde{c}_i$.} Substituting Eq.~\ref{eq:periodic_cov} into the double-integral form of the conditional variance,
% %
% \begin{align}
% 	\text{Var}[\tilde{c}_i \mid b_i]
% 	&= \int_0^\infty\!\!\int_0^\infty b_i(\tau)\,b_i(\tau')\,\frac{A^2}{2}\cos\bigl(\omega(\tau - \tau')\bigr)\,d\tau\,d\tau' \notag\\
% 	&= \frac{A^2}{2}\,\text{Re}\!\left[\left(\int_0^\infty b_i(\tau)\,e^{i\omega\tau}\,d\tau\right)\!\left(\int_0^\infty b_i(\tau')\,e^{-i\omega\tau'}\,d\tau'\right)\right] \notag\\
% 	&= \frac{A^2}{2}\,\bigl|\hat{b}_i(\omega)\bigr|^2,
% 	\label{eq:var_cond}
% \end{align}
% %
% where $\hat{b}_i(\omega) = \int_0^\infty b_i(\tau)\,e^{-i\omega\tau}\,d\tau$ is the Fourier transform (equivalently, the characteristic function at $\omega$) of the biological profile.

% \paragraph{Evaluating $|\hat{b}_i(\omega)|^2$.} For the time-shifted Gamma model, $b_i(\tau) = f_\epsilon(\tau - l_i)$ with $f_\epsilon$ a $\text{Gamma}(\psi\alpha, \beta)$ density and random latent period $l_i$. A time shift multiplies the Fourier transform by a unit-modulus phase, so
% %
% \begin{equation}
% 	\hat{b}_i(\omega) = e^{-i\omega l_i}\,\hat{f}_\epsilon(\omega), \qquad \bigl|\hat{b}_i(\omega)\bigr|^2 = \bigl|\hat{f}_\epsilon(\omega)\bigr|^2,
% \end{equation}
% %
% with the $l_i$-dependence dropping out entirely. The characteristic function of $\text{Gamma}(\psi\alpha, \beta)$ is $\hat{f}_\epsilon(\omega) = (1 - i\omega/\beta)^{-\psi\alpha}$, giving
% %
% \begin{equation}
% 	\bigl|\hat{f}_\epsilon(\omega)\bigr|^2 = \left(1 + \frac{\omega^2}{\beta^2}\right)^{-\psi\alpha}.
% 	\label{eq:cf_gamma}
% \end{equation}
% %
% Because Eq.~\ref{eq:cf_gamma} is deterministic, the outer expectation over $b_i$ in the law of total variance is trivial:
% %
% \begin{equation}
% 	\text{Var}[\tilde{c}_i] = \frac{A^2}{2}\left(1 + \frac{\omega^2}{\beta^2}\right)^{-\psi\alpha}.
% 	\label{eq:var_tildec}
% \end{equation}

% \paragraph{Closed-form expression for $k$.} Combining Eqs.~\ref{eq:k_bridge} and~\ref{eq:var_tildec},
% %
% \begin{equation}
% 	\boxed{\;
% 	k = \frac{2}{A^2}\left(1 + \frac{\omega^2}{\beta^2}\right)^{\psi\alpha}.
% 	\;}
% 	\label{eq:k_periodic}
% \end{equation}
% %
% The expression has two factors with transparent interpretations. The prefactor $2/A^2$ is the reciprocal variance of a uniformly-phased sinusoid with amplitude $A$, and it sets the overall scale: larger contact-oscillation amplitude $\Rightarrow$ smaller $k$ $\Rightarrow$ more overdispersion. The bracketed factor $(1 + \omega^2/\beta^2)^{\psi\alpha}$ is the reciprocal of the squared characteristic function of $b_i$ at the contact frequency $\omega$, and it quantifies how much the biological profile's temporal window fails to average out the contact oscillation. This factor increases with $\psi$ for any $\omega > 0$, so narrower profiles always produce more overdispersion in the presence of any periodic contact variation.

% \paragraph{Limiting cases.}
% \begin{itemize}
% 	\item \textbf{Spike profile, $\psi \to 0$:} the bracketed factor becomes $1$, leaving $k = 2/A^2$. A spike profile samples $c_i$ at a single instant, so the conditional variance of $\tilde{c}_i$ equals $\text{Var}[c_i] = A^2/2$ regardless of $\omega$ or $\beta$.
% 	\item \textbf{Smooth profile, $\psi \to 1$:} $k = (2/A^2)(1 + \omega^2/\beta^2)^{\alpha}$, strictly larger than the spike value whenever $\omega > 0$.
% 	\item \textbf{Slow oscillation, $\omega/\beta \to 0$:} the contact process looks locally constant over the infection window, so $\tilde{c}_i \approx c_i(t_i)$ for any $\psi$, and $k \to 2/A^2$ independent of $\psi$.
% 	\item \textbf{Fast oscillation, $\omega/\beta \to \infty$:} the contact process averages out over any infection window with $\psi > 0$, driving $\text{Var}[\tilde{c}_i] \to 0$ and $k \to \infty$ (Poisson).
% 	\item \textbf{No contact variation, $A \to 0$:} $k \to \infty$ trivially.
% \end{itemize}
% %
% These limits are consistent with the intuition that overdispersion arises from incomplete averaging of contact variation by the biological profile window, and that the degree of averaging is controlled by both $\psi\alpha$ (profile width) and $\omega/\beta$ (ratio of profile timescale to contact timescale).

% \subsubsection{Renewal equation model}

% Early work by Hudson, and later by Kermack and McKendrick \citep{kermack_contribution_1927}, modeled the transmission of a permanently immunizing infection using an integral equation. In modern notation \citep{breda_formulation_2012}, the force of infection is
% %
% \begin{equation} 
% F(t) = \int_0^\infty F(t - \tau) S(t - \tau) A(\tau) d\tau 
% \label{eq:foi}
% \end{equation}
% %
% where $F(t)$ is the force of infection at time $t$, $S(t)$ is the density of susceptible individuals in the population at time $t$, and $A(\tau)$ is the infectiousness profile describing the expected contribution to the force of infection from an individual who was infected $\tau$ time units ago. The incidence of disease is the product of the force of infection and the susceptible density (\ie, $-\frac{dS}{dt} = F(t) S(t)$); thus, the infectiousness profile $A(\tau)$ is the fundamental object that determines how an epidemic unfolds. 

% \clearpage

% \subsubsection*{Bias in empirical growth rate estimation}

% The early-epidemic growth rate is estimated by fitting a Poisson generalised linear model (GLM) with log link and identity offset to daily incidence counts within a cumulative-case window $[Z_{\min}, Z_{\max}]$, where $Z_{\min} = 100$ and $Z_{\max} = 500$ in our simulations. The fitted slope estimates the exponential growth rate $\hat{r}$, which we compare to the theoretical Euler--Lotka rate $r = \beta(R_0^{1/\alpha} - 1)$. Three sources of downward bias in $\hat{r}$ relative to $r$ are present; we describe each in turn.

% \paragraph{Susceptible depletion (dominant source).} The Euler--Lotka growth rate $r$ assumes an effectively infinite susceptible pool. In a finite population of size $N$, the effective reproduction number at cumulative incidence $Z$ is $R_{\text{eff}}(Z) = R_0 (1 - Z/N)$, and the corresponding instantaneous growth rate is
% %
% \begin{equation}
% 	r_{\text{eff}}(Z) = \beta \Bigl[\bigl(R_0 (1 - Z/N)\bigr)^{1/\alpha} - 1\Bigr].
% \end{equation}
% %
% The Poisson GLM fits an average slope over the estimation window. Evaluating at the midpoint $Z_{\text{mid}} = (Z_{\min} + Z_{\max})/2$ and expanding for small $f = Z_{\text{mid}} / N$ gives the approximate relative bias:
% %
% \begin{equation}
% 	\frac{\hat{r} - r}{r}
% 	\approx -\frac{R_0^{1/\alpha} \cdot f}{\alpha \bigl(R_0^{1/\alpha} - 1\bigr)}.
% 	\label{eq:saturation_bias}
% \end{equation}
% %
% With $Z_{\min} = 100$, $Z_{\max} = 500$, and $N = 10{,}000$, the midpoint fraction is $f = 0.03$. Evaluating Eq.~\eqref{eq:saturation_bias} for each pathogen:

% \begin{center}
% \begin{tabular}{lccc}
% 	\hline\hline
% 	Pathogen & $\alpha$ & $R_0^{1/\alpha}$ & Predicted bias \\
% 	\hline
% 	Influenza & 2.32 & 1.35 & $-5.0\%$ \\
% 	Omicron   & 2.39 & 2.12 & $-2.4\%$ \\
% 	Measles   & 11.36 & 1.24 & $-1.3\%$ \\
% 	\hline\hline
% \end{tabular}
% \end{center}

% \noindent These predictions agree well with empirical simulation results (influenza: $-5.7\%$; omicron: $-3.5\%$; measles: $-0.1\%$). The quantity $R_0^{1/\alpha}/[\alpha(R_0^{1/\alpha} - 1)]$ in Eq.~\eqref{eq:saturation_bias} is the elasticity of $r$ with respect to $R_0$: it measures how sensitively the growth rate responds to changes in the effective reproduction number. This elasticity is largest when $R_0$ is close to the epidemic threshold ($R_0^{1/\alpha} \to 1$), because the growth rate $r = \beta(R_0^{1/\alpha} - 1)$ vanishes at $R_0 = 1$ and is steeply increasing just above it. Influenza ($R_0 = 2$) has the highest elasticity (1.67), so even a 3\% depletion of susceptibles produces a $-5\%$ reduction in $r$. Measles ($R_0 = 12$) has the lowest elasticity (0.45), both because it is far from threshold and because its large $\alpha$ ensures that $R_0^{1/\alpha}$ is close to 1 despite the high $R_0$. The bias can be reduced by using a narrower window (smaller $Z_{\max}$) or a larger population $N$, at the cost of increased estimator variance or computational expense.

% \paragraph{Boundary-day truncation (corrected).} Because the estimation window is defined by cumulative case counts rather than calendar time, the first and last calendar days of each simulation's window are partial: the first day contains only infections occurring after the cumulative count crosses $Z_{\min}$, and the last day contains only infections before the count reaches $Z_{\max}$. Daily incidence counts on these boundary days are therefore systematically lower than the true daily rate. In a Poisson GLM, the high-count final day receives disproportionate weight (the Poisson likelihood weights observations inversely by their variance, which equals the mean), so even modest truncation of the last day's count pulls the fitted slope downward. In an unweighted log-linear fit, the effect is smaller but still present. We verified this bias empirically using a non-homogeneous Poisson process (NHPP) simulator with known rate $\lambda(t) = \lambda_0 e^{rt}$: with boundary days included, the Poisson GLM underestimated $r$ by approximately $34\%$; dropping the first and last calendar days from the fit eliminated the bias (residual $< 0.5\%$). All growth rate estimates reported in this paper use the corrected estimator with boundary days removed.

% \paragraph{Generational cohort structure (affects spike profiles).} For highly punctuated profiles ($\psi \to 0$), secondary infections arrive in discrete cohorts separated by approximately one mean generation time $\bar{g} = \alpha/\beta$. If the estimation window spans fewer than approximately two generation times in calendar time, the daily incidence series may capture the rise and fall of a single generational cohort rather than the average exponential trend. This produces noisy slope estimates and occasional negative values. The effect is most pronounced for measles at $\psi = 0$: the generation time is $\bar{g} = 12.2$ days while the $[100, 500]$ window spans only $\approx \log(5)/r \approx 7$ days, less than one full generation. For influenza ($\bar{g} = 3.2$ days, window $\approx 6.4$ days) and omicron ($\bar{g} = 7.1$ days, window $\approx 4.3$ days), the window covers at least one full generation and the cohort effect is attenuated. This source of bias is inherent to the spike model and diminishes as $\psi$ increases: for $\psi = 1$, offspring times are independent and daily incidence follows a smooth exponential trend even over short windows.

% \paragraph{Summary.} The dominant source of downward bias is susceptible depletion, which is quantitatively predicted by Eq.~\eqref{eq:saturation_bias} and is independent of the punctuation parameter $\psi$. After correcting for boundary-day truncation, the residual bias is consistent with the depletion prediction for all three pathogens. The theoretical growth rate $r$ from the Euler--Lotka equation remains the correct characterisation of the early epidemic; the empirical estimates are biased downward because they necessarily operate in a regime where finite-population effects have begun to reduce the effective reproduction number below $R_0$.

% \clearpage

% \subsubsection*{Variance of the growth rate estimator and dependence on $\psi$}

% In simulations, the standard deviation of the empirical growth rate $\hat{r}$ increases monotonically as $\psi \to 0$ (Table~\ref{tab:growthrate_sd}). Here we provide an analytic explanation for this pattern, based on the overdispersion of daily incidence counts induced by offspring synchronisation.

% \begin{table}[h]
% \centering
% \begin{tabular}{llccc}
% 	\hline\hline
% 	Pathogen & $\psi$ & $n$ (established) & Mean $\hat{r}$ & SD $\hat{r}$ \\
% 	\hline
% 	Influenza & 0   & 792  & 0.195 & 0.066 \\
% 	          & 0.5 & 788  & 0.197 & 0.058 \\
% 	          & 1   & 803  & 0.197 & 0.052 \\[4pt]
% 	Omicron   & 0   & 997  & 0.281 & 0.127 \\
% 	          & 0.5 & 997  & 0.275 & 0.086 \\
% 	          & 1   & 999  & 0.279 & 0.079 \\[4pt]
% 	Measles   & 0   & 1000 & 0.227 & 0.128 \\
% 	          & 0.5 & 1000 & 0.200 & 0.054 \\
% 	          & 1   & 1000 & 0.191 & 0.043 \\
% 	\hline\hline
% \end{tabular}
% \caption{Mean and standard deviation of the empirical growth rate estimator $\hat{r}$ across established epidemics, by pathogen and punctuation parameter $\psi$. Based on 1000 simulations per pathogen with $N = 10{,}000$, estimation window $[Z_{\min}, Z_{\max}] = [100, 500]$.}
% \label{tab:growthrate_sd}
% \end{table}

% \paragraph{Per-infector index of dispersion.} Consider a single infector who produces $\chi \sim \text{Poisson}(R_0)$ offspring. Let $q_t = P(\tau \in [t, t+1))$ denote the probability that a generation interval $\tau \sim g(\tau)$ falls in the calendar day $[t, t+1)$, and let $N_i(t)$ denote the number of this infector's offspring whose infection times fall on day~$t$.

% In the \textbf{smooth model} ($\psi = 1$), offspring arrive at independent times $\tau_j \sim g(\tau)$. Conditional on $\chi$, the count on day $t$ is $\text{Binomial}(\chi, q_t)$. Marginalising over $\chi \sim \text{Poisson}(R_0)$ gives $N_i(t) \sim \text{Poisson}(R_0 q_t)$, with index of dispersion
% %
% \begin{equation}
% 	\text{ID}_{\text{smooth}} = \frac{\text{Var}[N_i(t)]}{E[N_i(t)]} = 1.
% \end{equation}

% In the \textbf{spike model} ($\psi = 0$), all $\chi$ offspring land at the same instant $\tau^* \sim g(\tau)$. The count on day $t$ is either $0$ (with probability $1 - q_t$) or $\chi$ (with probability $q_t$). Computing the first two moments:
% %
% \begin{align}
% 	E[N_i(t)] &= R_0 \, q_t, \\[4pt]
% 	E[N_i(t)^2] &= E[\chi^2] \cdot q_t = (R_0 + R_0^2) \, q_t,
% \end{align}
% %
% so $\text{Var}[N_i(t)] = R_0 \, q_t + R_0^2 \, q_t(1 - q_t)$ and
% %
% \begin{equation}
% 	\text{ID}_{\text{spike}} = 1 + R_0(1 - q_t).
% 	\label{eq:id_spike}
% \end{equation}
% %
% Since $q_t \ll 1$ for any single day (the generation interval spans multiple days), $\text{ID}_{\text{spike}} \approx 1 + R_0$.

% \paragraph{General $\psi$.} For intermediate $\psi$, the offspring from infector $i$ are born at times $\tau_j = l_i + \epsilon_j$, where $l_i \sim \text{Gamma}((1-\psi)\alpha, \beta)$ is the shared latent period and $\epsilon_j \sim \text{Gamma}(\psi\alpha, \beta)$ is independent jitter. Conditional on $l_i$, the $\epsilon_j$ are independent, so the count on day $t$ given $l_i$ is $\text{Poisson}(R_0 \cdot p(l_i, t))$ where $p(l, t) = P(\epsilon \in [t - l, t - l + 1))$. Marginalising over $l_i$ using the law of total variance:
% %
% \begin{equation}
% 	\text{Var}[N_i(t)] = R_0 \, q_t + R_0^2 \, \text{Var}_{l_i}\bigl[p(l_i, t)\bigr],
% \end{equation}
% %
% giving index of dispersion
% %
% \begin{equation}_{}
% 	\text{ID}(\psi) = 1 + R_0 \cdot \frac{\text{Var}_{l_i}[p(l_i, t)]}{q_t}.
% 	\label{eq:id_psi}
% \end{equation}
% %
% The quantity $\text{Var}_{l_i}[p(l_i, t)] / q_t$ interpolates monotonically from $1 - q_t \approx 1$ at $\psi = 0$ (where $p(l_i, t) \approx \mathbf{1}(l_i \in [t, t+1))$ is nearly Bernoulli) to $0$ at $\psi = 1$ (where $l_i \approx 0$ and $p$ is deterministic). Thus the per-infector index of dispersion decreases monotonically from $1 + R_0$ to $1$ as $\psi$ increases from $0$ to $1$.

% \paragraph{Aggregation.} The total daily incidence on day $t$ is $\Delta Z(t) = \sum_i N_i(t)$, summing over all previously infected individuals. Conditional on the infection times $\{t_i\}$, the contributions $N_i$ from distinct infectors are independent (each infector's offspring process is independent of every other's). Therefore, by the additivity of variances:
% %
% \begin{equation}
% 	\text{Var}[\Delta Z(t) \mid \{t_i\}] = \sum_i \text{Var}[N_i(t)] = \sum_i \bigl[R_0 \, q_{t-t_i} + R_0^2 \, \text{Var}_l[p(l, t-t_i)]\bigr].
% \end{equation}
% %
% Since $E[\Delta Z(t) \mid \{t_i\}] = \sum_i R_0 \, q_{t-t_i}$, the aggregate index of dispersion is
% %
% \begin{equation}
% 	\phi(\psi) = 1 + R_0 \cdot \frac{\sum_i \text{Var}_l[p(l, t - t_i)]}{\sum_i q_{t-t_i}},
% \end{equation}
% %
% a weighted average of the per-infector overdispersion across all contributing infectors. The overdispersion does not wash out with aggregation because each infector's contribution is independently overdispersed by the same mechanism. The total daily incidence is therefore quasi-Poisson with dispersion parameter $\phi(\psi)$ that interpolates from $\phi(0) \approx 1 + R_0$ to $\phi(1) = 1$.

% \paragraph{Implication for growth rate estimation.} Because the growth rate $\hat{r}$ is estimated by regressing log daily incidence on time, the variance of daily incidence directly inflates the variance of $\hat{r}$. Since the aggregate dispersion parameter $\phi(\psi)$ increases monotonically from $1$ at $\psi = 1$ to approximately $1 + R_0$ at $\psi = 0$, the growth rate estimator is correspondingly more variable for more punctuated profiles. This prediction is confirmed by the simulation results in Table~\ref{tab:growthrate_sd}: for all three pathogens, $\text{SD}(\hat{r})$ increases as $\psi$ decreases from $1$ to $0$, with the effect most pronounced for measles ($R_0 = 12$), where the dispersion ratio $\phi(0)/\phi(1) \approx 1 + R_0 = 13$ is largest.

\clearpage 

\subsubsection*{Calculating overdispersion}

{\bf \boldmath Deriving an expression for the overdispersion parameter $k$.} We begin with the standard assumption that the number of offspring $Z_i$ produced by an index case $i$ is Poisson-distributed: 
\begin{equation}
	Z_i \sim \text{Poisson}(\nu_i)
\end{equation}
When $\nu_i$ is uniform across people ($\nu_i \equiv R_0$), then the $Z_i$ are truly Poisson-distributed. Overdispersion occurs when $\nu_i$ varies across individuals. To capture this overdispersion, the standard approach is to describe $Z_i$ using a Negative Binomial distribution: 
\begin{equation}
	Z_i \sim \text{NegBin}(R_0, k)
\end{equation}
which has expectation $R_0$ and variance $R_0 (1 + R_0 / k)$. Thus, $k$ captures overdispersion, with $k \rightarrow \infty$ yielding Poisson-distributed $Z_i$ and $k \rightarrow 0$ yielding increasingly heavier Negative Binomial tails (\ie, more overdispersion). The Negative Binomial expression is exact when $\nu_i$ are Gamma-distributed; otherwise, it is an approximation to the distribution of $Z_i$. Either way, the parameter $k$ can be obtained {\it via} moment matching by re-arranging the variance expression: 
\begin{equation}
	k = \frac{R_0^2}{\text{Var}[Z_i] - R_0}
\end{equation}
Furthermore, since
\begin{equation}
	\text{Var}[Z_i] = \text{Var}[E[Z_i | \nu_i]] + E[\text{Var}[Z_i | \nu_i]] = \text{Var}[\nu_i] + E[\nu_i] = \text{Var}[\nu_i] + R_0
\end{equation}
the expression for $k$ simplifies to 
\begin{equation}
	k = \frac{R_0^2}{\text{Var}[\nu_i]}
	\label{eq:k_expression}
\end{equation}

Next, if we assume everyone's inherent biological infectiousness is the same ($b_i \equiv R_0$), we have 
\begin{equation}
	\nu_i = R_0 \int_0^\infty g_i(\tau) c_i(t_i + \tau) d\tau := R_0 \tilde{c}_i
\end{equation}
Plugging in to \ref{eq:k_expression} gives 
\begin{equation}
	\boxed{k = \frac{1}{\text{Var}[\tilde{c}_i]}}
\end{equation}
Thus, the overdispersion $k$ depends only on the variance of $\tilde{c}_i = \int_0^\infty g_i(\tau) c_i(t_i + \tau) d\tau$. 

\noindent {\bf \boldmath Deriving an expression for $\text{Var}[\tilde{c}_i]$.} The random variable $\tilde{c}_i$ inherits randomness from both $g_i(\tau)$ and $c_i(t_i+\tau)$. To manage this, we again apply the Law of Total Variance: 
\begin{align}
	\text{Var}[\tilde{c}_i] &= \text{Var}[E[\tilde{c}_i | g_i ]] + E[\text{Var}[\tilde{c}_i | g_i ]] \\ 
	&= \text{Var}[1] + E[\text{Var}[\tilde{c}_i | g_i ]] \\ 
	&= E[\text{Var}[\tilde{c}_i | g_i ]]
\end{align}
Thus, we are interested in 
\begin{equation}
\tilde{c}_i | g_i = \int_0^\infty g_i(\tau) c_i(t_i+\tau) d\tau 
\end{equation}
where $g_i$ is known (\ie, conditioned upon). This has the form of a classic problem in signal processing: $c_i$ is a wide-sense stationary stochastic process (an input signal); $g_i$ is an impulse response; and $\tilde{c}_i | g_i$ is the output signal evaluated at time $t_i$. Together, this comprises a linear system. Following standard theory, we can express $\text{Var}[\tilde{c}_i | g_i]$ in terms of the autocorrelation $R(\tau)$ of the signal $\tilde{c}_i | g_i$ at time lag $\tau$. Specifically, 
\begin{align}
	\text{Var}[\tilde{c}_i | g_i] &= E[(\tilde{c}_i | g_i)^2]  - E[\tilde{c}_i | g_i]^2  \\ 
	&= 	R_{\tilde{c}_i | g_i}(0) - E[\tilde{c}_i | g_i]^2 \\
	&= 	R_{\tilde{c}_i | g_i}(0) - 1 
\end{align}

\noindent {\bf \boldmath Deriving an expression for $R_{\tilde{c}_i | g_i}(0)$.} Deriving an expression for the second moment of the output process, $R_{\tilde{c}_i | g_i}(0)$, is easiest in frequency space. We can introduce $S_\bullet$ as the Fourier transform of $R_\bullet$ and $\hat{g}_i$ as the Fourier transform of $g_i$, \ie 
\begin{equation}
S_\bullet(\omega) = \int_{-\infty}^{\infty} e^{-i \omega \tau} R_\bullet(\tau) d\tau 
\qquad \text{ and } \qquad 
\hat{g}_i(\omega) = \int_{-\infty}^{\infty} e^{-i \omega \tau} g_i(\tau) d\tau
\end{equation}
By the inverse Fourier transform, 
\begin{equation}
R_\bullet(\tau) = \frac{1}{2\pi} \int_{-\infty}^{\infty} S_\bullet(\omega) e^{i \omega \tau} d\omega
\end{equation}
and thus 
\begin{equation}
R_\bullet(0) = \frac{1}{2\pi} \int_{-\infty}^{\infty} S_\bullet(\omega) d\omega
\end{equation}
We thus seek an expression for $S_{\tilde{c}_i | g_i}(\omega)$, which we can obtain from the fundamental theorem for linear systems: 
\begin{equation}
	S_{\tilde{c}_i | g_i}(\omega) = |\hat{g}_i(\omega)|^2 S_{c_i}(\omega)
\end{equation}

\noindent {\bf \boldmath Deriving $\hat{g}_i(\omega).$} We begin by examining $\hat{g}_i(\omega)$. Under the Gamma time-shift model, $g_i(\tau) = f_\epsilon(\tau - l_i)$, where $f_\epsilon$ is the $\text{Gamma}(\psi \alpha, \beta)$ density and $l_i \sim \text{Gamma}((1-\psi) \alpha, \beta)$. Taking the Fourier transform, 
\begin{align}
	\hat{g}_i(\omega) &= \int_{-\infty}^{\infty} g_i(\tau) e^{-i \omega \tau} d\tau \\ 
	&=	\int_{l_i}^{\infty} f_\epsilon(\tau - l_i) e^{-i \omega \tau}  d\tau\\ 
	&= \int_0^\infty f_\epsilon(\tau) e^{-i \omega (\tau + l_i)} d\tau \\
	&=  e^{-i \omega l_i} \int_0^\infty f_\epsilon(\tau) e^{-i \omega \tau} d\tau \\ 
	&= e^{-i \omega l_i} \int_0^\infty \frac{\beta^{\psi \alpha}}{\Gamma(\psi \alpha)} \tau^{\psi \alpha - 1} e^{-\beta \tau} e^{-i\omega\tau} d\tau \\
&= e^{-i \omega l_i} \Bigl(  \frac{\beta}{\beta + i \omega}\Bigr)^{\psi \alpha}
\end{align}
Taking the squared modulus gives 
\begin{equation}
 |\hat{g}_i(\omega)|^2 = \Bigl(1 + \frac{\omega^2}{\beta^2}\Bigr)^{-\psi \alpha}
\end{equation}
This is the power transfer function; it is constant across individuals and does not depend on $l_i$. 

\noindent {\bf \boldmath Deriving $S_{c_i}(\omega)$ for the periodic contact model.} For the periodic contact model
\begin{equation}
c_i(t) = 1 - c_{amp} \cos(\omega_0 t)
\end{equation}
the autocorrelation is 
\begin{align}
R_{c_i}(\tau)  &= E[(1 - c_{amp} \cos(\omega_0 t)) (1 - c_{amp} \cos(\omega_0 (t + \tau)))] \\ 
&= E[1] - c_{amp}E[\cos(\omega_0 t)] - c_{amp}E[\cos(\omega_0 (t + \tau))] + c_{amp}^2 E[\cos(\omega_0 t) \cos(\omega_0 (t + \tau))] \\ 
&= 1 + \frac{c_{amp}^2}{2} \cos(\omega_0 \tau) 
\end{align}
Taking the Fourier transform, 
\begin{align}
	S_{c_i}(\omega) &= \int_{-\infty}^{\infty} \Bigl[ 1 + \frac{c_{amp}^2}{2} \cos(\omega_0\tau) \Bigr] e^{-i\omega\tau}d\tau \\ 
	&= 2 \pi \delta(\omega) + \frac{\pi c_{amp}^2}{2}[\delta(\omega - \omega_0) + \delta(\omega + \omega_0)] 
\end{align}

\noindent {\bf \boldmath Constructing $k$ for the periodic contact model.} Putting everything together, we have: 
\begin{align}
\text{Var}[\tilde{c}_i | g_i] &= R_{\tilde{c}_i | g_i}(0) - 1 \\
&= \frac{1}{2\pi} \int_{-\infty}^{\infty} S_{\tilde{c}_i | g_i}(\omega) d\omega  - 1 \\
&= \frac{1}{2\pi} \int_{-\infty}^\infty \Bigl(1 + \frac{\omega^2}{\beta^2}\Bigr)^{-\psi \alpha} 
\Bigl[ 2 \pi \delta(\omega) + \frac{\pi c_{amp}^2}{2} [\delta(\omega - \omega_0) + \delta(\omega + \omega_0)] \Bigr] d \omega  - 1\\ 
&= 1 + \frac{c_{amp}^2}{2} \Bigl(1 + \frac{\omega_0^2}{\beta^2}\Bigr)^{-\psi \alpha} - 1 \\
&= \frac{c_{amp}^2}{2} \Bigl(1 + \frac{\omega_0^2}{\beta^2}\Bigr)^{-\psi \alpha}
\end{align}
Thus, 
\begin{equation}
	\boxed{k = \frac{2}{c_{amp}^2} \Bigl(1 + \frac{\omega_0^2}{\beta^2}\Bigr)^{\psi \alpha}}
\end{equation}

\noindent {\bf \boldmath Deriving $S_{c_i}(\omega)$ for the Gamma-Poisson contact model.} Let $X \sim \text{Gamma}(k_c, k_c)$ be a contact-level draw from the Gamma-Poisson contact model. Then, the autocorrelation 
\begin{equation}
	R_{c_i}(\tau) = E[c_i(t) c_i(t + \tau)]
\end{equation}
is equal to $E[X^2] = \text{Var}[X] + E[X]^2 = \frac{1}{k_c} + 1$ with probability $e^{-\lambda |\tau|}$ (\ie, if the contact levels have not switched between times $t$ and $t + \tau$), and is equal to $E[X]^2 = 1$ otherwise. That is: 
\begin{align}
	R_{c_i}(\tau) &= (\frac{1}{k_c} +1) e^{-\lambda |\tau|} + (1)(1 - e^{-\lambda |\tau|}) \\ 
	&= \frac{1}{k_c} e^{-\lambda |\tau|} + 1
\end{align}
Thus: 
\begin{align}
S_{c_i}(\omega) &= \int_{-\infty}^{\infty} [\frac{1}{k_c}e^{-\lambda |\tau| } + 1] e^{-i \omega \tau} d\tau \\ 
&= \frac{1}{k_c} \int_{-\infty}^{0} e^{ \lambda \tau} e^{- i \omega \tau} d\tau + 
\frac{1}{k_c} \int_{0}^{\infty} e^{- \lambda \tau } e^{- i \omega \tau} d\tau + 2\pi \delta(\omega) \\ 
&= \frac{1}{k_c} \Bigl[ \frac{1}{\lambda - i \omega} + \frac{1}{\lambda + i \omega} \Bigr]  + 2\pi \delta(\omega) \\
&= \frac{2 \lambda}{k_c (\lambda^2 + \omega^2)} + 2\pi \delta(\omega) 
\end{align}

\noindent {\bf \boldmath Constructing $k$ for the Gamma-Poisson contact model.} Putting everything together, we have: 
\begin{align}
\text{Var}[\tilde{c}_i | g_i] &= R_{\tilde{c}_i | g_i}(0) - 1 \\
&= \frac{1}{2\pi} \int_{-\infty}^{\infty} 
\Bigl[\frac{2 \lambda}{k_c (\lambda^2 + \omega^2)} + 2\pi \delta(\omega)\Bigr]
\Bigl( 1 + \frac{\omega^2}{\beta^2}\Bigr )^{-\psi \alpha} d\omega - 1 \\
&= \frac{1}{2\pi} \int_{-\infty}^{\infty} 
\Bigl[\frac{2 \lambda}{k_c (\lambda^2 + \omega^2)} \Bigr] \Bigl( 1 + \frac{\omega^2}{\beta^2}\Bigr )^{-\psi \alpha} d\omega \\
&= \frac{1}{k_c}\int_{-\infty}^{\infty} \frac{\lambda}{\pi (\lambda^2 + \omega^2)} \Bigl( 1 + \frac{\omega^2}{\beta^2}\Bigr )^{-\psi \alpha} d\omega
\end{align}
and $k$ is the reciprocal of this integral. The integral can be solved in closed form using special functions, but its final form is not especially illuminating; in practice, we evaluate this expression by numerically solving the integral. 

\clearpage

\subsection*{Supplementary Figures}

\subsection*{Supplementary Tables}


